% Revised working copy of F023, matched to F024.
% M08: draw y=8x-10, not a vertical line.
% M15: monotonicity is per connected interval; -3 is not an inflection point.
\documentclass[11pt,a4paper]{article}
% Shared print layout. The web's provisional design and MathJax are independent.
\usepackage[utf8]{inputenc}
\usepackage[T1]{fontenc}
\usepackage[spanish,es-nodecimaldot]{babel}
\usepackage{lmodern}
\usepackage{amsmath,amssymb}
\usepackage{graphicx}
\usepackage{geometry}
\geometry{a4paper,left=20mm,right=20mm,top=25mm,bottom=22mm,headheight=16pt,headsep=8mm}
\usepackage{fancyhdr,lastpage}
\usepackage{enumitem}
\usepackage{xcolor}
\usepackage{pgfplots}
\pgfplotsset{compat=1.18}
\usepgfplotslibrary{fillbetween}
\usepackage[hidelinks,pdfencoding=auto]{hyperref}
\hypersetup{pdflang=es-ES,pdfauthor={},pdfsubject={Matemáticas Aplicadas a las Ciencias Sociales II}}
\pagestyle{fancy}
\fancyhf{}
\fancyhead[L]{\small Matemáticas Aplicadas a las Ciencias Sociales II}
\fancyhead[R]{\small 2.\textsuperscript{o} Bachillerato}
\fancyfoot[L]{\footnotesize Recuperación T1 | Edición revisada 1.0}
\fancyfoot[R]{\footnotesize Página \thepage\ de \pageref{LastPage}}
\renewcommand{\headrulewidth}{0.4pt}
\setlength{\parindent}{0pt}
\setlength{\parskip}{0.5em}
\setlist{itemsep=0.3em,topsep=0.3em}
\newcommand{\apartado}[1]{\subsection*{#1}}
\pgfplotsset{
  revision/.style={
    width=0.94\linewidth,height=7.5cm,
    axis lines=middle,grid=major,grid style={gray!18},
    tick label style={font=\small},label style={font=\small},
    xlabel={$x$},ylabel={$y$},samples=160,
    legend style={at={(0.5,-0.16)},anchor=north,draw=none,font=\small},
    clip=true
  }
}

\hypersetup{pdftitle={Recuperación T1 - Soluciones revisadas}}
\begin{document}
{\Large\bfseries Recuperación T1: soluciones}

\section*{Ejercicio 1(a): continuidad y derivabilidad}
\textbf{Datos del enunciado.}
\[
f(x)=
\begin{cases}
ax-3x^2,& x\leq1,\\
2x^2+b,& x>1.
\end{cases}
\]
Para que \(f\) sea derivable en \(x=1\), primero debe ser continua allí.
Después comprobaremos la igualdad de sus derivadas laterales.

\apartado{1. Continuidad en el punto de unión}
El valor en el punto y los límites laterales son
\[
f(1)=a-3,\qquad
\lim_{x\to1^-}f(x)=a-3,\qquad
\lim_{x\to1^+}f(x)=2+b.
\]
Por tanto, la condición de continuidad es
\[
a-3=2+b \quad\Longleftrightarrow\quad a-b=5.
\]

\apartado{2. Derivadas laterales}
En el interior de cada rama,
\[
f'(x)=
\begin{cases}
a-6x,& x<1,\\
4x,& x>1.
\end{cases}
\]
Bajo la condición de continuidad, las derivadas laterales en la unión son
\[
f'_-(1)=a-6,\qquad f'_+(1)=4.
\]
Igualándolas se obtiene \(a-6=4\), es decir, \(a=10\). Al sustituir en
\(a-b=5\), resulta \(b=5\).

\[
\boxed{a=10,\qquad b=5}
\]
\textbf{Comprobación:} ambas ramas dan \(7\) en la unión y las dos
derivadas laterales valen \(4\). Con estos parámetros existe \(f'(1)=4\).
Las ramas son polinomios, así que no hay otros puntos problemáticos.

\clearpage
\section*{Ejercicio 1(b): recta tangente en \(x=2\)}
Este apartado fija \textbf{otros parámetros}: \(a=3\) y \(b=-2\).
No se utilizan los valores obtenidos en el apartado anterior.
Como \(2>1\), alrededor del punto de tangencia rige
\[
f(x)=2x^2-2,\qquad f'(x)=4x.
\]
El punto y la pendiente son
\[
P=(2,f(2))=(2,6),\qquad m=f'(2)=8.
\]
La ecuación punto-pendiente proporciona
\[
y-6=8(x-2)
\quad\Longrightarrow\quad
\boxed{y=8x-10}.
\]
\textbf{Comprobación:} al sustituir \(x=2\) la recta da \(y=6\), y su
pendiente es \(8\), igual a \(f'(2)\). Una recta vertical \(x=2\) no
representa esta tangente.

\begin{center}
\begin{tikzpicture}
\begin{axis}[revision,xmin=-1.1,xmax=3.2,ymin=-8,ymax=17,
  xtick={-1,0,1,2,3},ytick={-5,0,5,10,15},height=8.1cm]
  \addplot[black,thick,domain=-1:1]{3*x-3*x^2};
  \addlegendentry{\(3x-3x^2\), para \(x\leq1\)}
  \addplot[blue!70!black,thick,domain=1:3]{2*x^2-2};
  \addlegendentry{\(2x^2-2\), para \(x>1\)}
  \addplot[black,thick,dashed,domain=0.5:3.2]{8*x-10};
  \addlegendentry{Tangente: \(y=8x-10\)}
  \addplot[only marks,mark=*,mark size=2pt,black] coordinates {(2,6)};
  \node[anchor=north west,font=\small] at (axis cs:2,6){\(P=(2,6)\)};
\end{axis}
\end{tikzpicture}
\end{center}

\clearpage
\section*{Ejercicio 2(a): estudio de \(g(x)=1/(x+3)\)}
\textbf{Dominio y continuidad.}
El denominador se anula en \(-3\), de modo que
\[
D_g=\mathbb{R}\setminus\{-3\}.
\]
La función es continua en cada intervalo de su dominio. En \(x=-3\)
presenta una discontinuidad infinita.

\textbf{Asíntotas.} Los límites laterales son
\[
\lim_{x\to-3^-}g(x)=-\infty,\qquad
\lim_{x\to-3^+}g(x)=+\infty.
\]
La asíntota vertical es \(x=-3\). Además,
\(\lim_{x\to-\infty}g(x)=\lim_{x\to+\infty}g(x)=0\), por lo que la
asíntota horizontal es \(y=0\).

\textbf{Monotonía y curvatura.}
\[
g'(x)=-\frac{1}{(x+3)^2}<0,\qquad
g''(x)=\frac{2}{(x+3)^3}\quad(x\ne-3).
\]
\begin{center}
\begin{tabular}{c|c|c}
Intervalo & Monotonía & Curvatura\\ \hline
\((-\infty,-3)\) & Decreciente & Hacia abajo: \(g''<0\)\\
\((-3,+\infty)\) & Decreciente & Hacia arriba: \(g''>0\)
\end{tabular}
\end{center}
La disminución es \textbf{en cada intervalo}, no en todo el dominio
considerado conjuntamente: \(-4<-2\), pero \(g(-4)=-1<1=g(-2)\).
No hay extremos relativos. Tampoco hay punto de inflexión en \(-3\),
porque ese valor no pertenece al dominio. La curva no corta el eje \(x\)
y corta el eje \(y\) en \((0,1/3)\).

\begin{center}
\begin{tikzpicture}
\begin{axis}[revision,xmin=-7,xmax=3,ymin=-3,ymax=3,height=5.7cm,
  xtick={-6,-3,0,3},ytick={-2,-1,0,1,2},
  legend style={at={(0.5,-0.20)},anchor=north,draw=none,font=\small,legend columns=2}]
  \addplot[black,thick,domain=-7:-3.12]{1/(x+3)};
  \addlegendentry{\(g(x)=1/(x+3)\)}
  \addplot[black,thick,domain=-2.88:3,forget plot]{1/(x+3)};
  \addplot[black,dashed,thick] coordinates {(-3,-3)(-3,3)};
  \addlegendentry{\(x=-3\)}
  \addplot[black,dotted,thick] coordinates {(-7,0)(3,0)};
  \addlegendentry{\(y=0\)}
\end{axis}
\end{tikzpicture}
\end{center}

\clearpage
\section*{Ejercicio 2(b): área entre \(x=2\) y \(x=4\)}
En \([2,4]\), el denominador \(x+3\) es positivo. La gráfica está por
encima del eje de abscisas, así que el área coincide con la integral:
\[
A=\int_2^4\frac{1}{x+3}\,dx.
\]
Una primitiva en este intervalo es \(\ln|x+3|\). Aplicando la regla de
Barrow,
\begin{align*}
A&=\bigl[\ln|x+3|\bigr]_2^4\\
 &=\ln 7-\ln 5\\
 &=\ln\!\left(\frac75\right).
\end{align*}
\[
\boxed{A=\ln\!\left(\frac75\right)\ \text{unidades}^2
       \approx 0{,}33647\ \text{unidades}^2}
\]
\textbf{Comprobación del orden de magnitud.}
Como \(g\) decrece en \([2,4]\),
\[
\frac{2}{7}\leq A\leq\frac{2}{5}.
\]
El valor \(0{,}33647\) satisface estas cotas y es positivo.

\begin{center}
\begin{tikzpicture}
\begin{axis}[revision,xmin=1,xmax=5,ymin=0,ymax=0.28,height=7cm,
  xtick={1,2,3,4,5},ytick={0,0.1,0.2},
  legend style={at={(0.5,-0.18)},anchor=north,draw=none,font=\small}]
  \addplot[name path=upper,draw=none,domain=2:4,forget plot]{1/(x+3)};
  \path[name path=lower] (axis cs:2,0)--(axis cs:4,0);
  \addplot[draw=none,fill=gray!25,forget plot] fill between[of=upper and lower];
  \addplot[black,thick,domain=1:5]{1/(x+3)};
  \addlegendentry{\(g(x)=1/(x+3)\)}
  \addplot[black,dashed,forget plot] coordinates {(2,0)(2,0.2)};
  \addplot[black,dashed,forget plot] coordinates {(4,0)(4,0.142857142857)};
\end{axis}
\end{tikzpicture}
\end{center}
\end{document}
